The strong interaction, described by QCD, is one of the four fundamental forces and governs the behavior of quarks and gluons. Unlike electromagnetism, where the force is mediated by photons that do not interact with each other, the strong interaction is mediated by gluons, which are massless bosons, charged under the strong force, with charge which is referred to as color. Contrary to the electric charge, color charge has 6 unique elements, labeled as red, green, blue and their counterparts (i.e - anti-red). The strong interaction is responsible for binding the nucleons, which are comprised of thee generations of quarks (anti-quarks) and gluons together, such as proton (nominally uud) and neutrons ( nominally udd).

As a result of the symmetry and non-Abelian nature of QCD, there cannot be free particles with color charge. Instead, all observed particles are color-singlets. Contrary to electromagnetism, one of the proposed potential of the strong interaction, is linear with distance at large distances (Cornell potential \cite{brambilla_quark_1999}):
\begin{equation}\label{eq:qcd_potential}
V(r) = -\frac{4}{3}\frac{\alpha_s}{r} + \sigma r + Const \,,
\end{equation}

where $\alpha_s$ is the strong coupling and $\sigma$ is called the string tension, related to the flux tubes connecting color-charged partons. Due to this potential, as color charged particles drift apart, the potential energy increases, until eventually it is preferable to split the flux tube between the particles and generate new color charged particles in place, as illustrated in Fig. \ref{fig:confinement}.

\begin{figure}[h!] 
    \centering % Centers the image and caption
    \includegraphics[width=0.5\textwidth]{\main/illustrations/color_confinement.png}
    \caption{Illustration of a flux tube breaking, generating new quarks and anti-quarks in the process \cite{greensite_introduction_2020}.}
    \label{fig:confinement}
\end{figure}

\subsection{QCD Jet}\label{sec:qcdjet}
As discussed previously, quarks and gluons are bound by the strong interaction and will generate new color partons when pulled apart. Following that, free quarks and gluons in a highly boosted frame of reference will generate a shower of color partons through a bremsstrahlung cascade,  which will subsequently undergo hadronization (see Fig. \ref{fig:qcd_jet}) and will result in highly collimated clusters of hadrons, leptons and photons. We usually refer to the initial instigators of the cascade as the QCD jets \cite{ali_jets_2011}. Additional information regarding the measurable quantity known as jet will be explained in section \ref{sec:reco_jet}.


 \begin{figure}[h!] 
 	\centering % Centers the image and caption
 	\includegraphics[width=0.8\textwidth]{\main/illustrations/Ali_qcd_fragmantation.png}
 	\caption{(a) Quark fragmentation to hadrons induced by confinement forces; (b) quark/gluon cascades at high energies in QCD \cite{ali_jets_2011}.}
 	\label{fig:qcd_jet}
 \end{figure}

The process of hadronization is a topic of active research. Currently research into hadronization uses various models and simulations, as perturbative QCD is not applicable at these energies. 
